% !PRINTED 376
\weberpage{376}
%<<CONT
\weq{10}{S_n \geqq S' \geqq 0,}
in which, however, the equality sign can occur in the lower limit only for $m = 1$, and in the upper limit only for $m = 2$ and $b_1 = 1$, hence in general only when $S_n = 1$; and this can be avoided, by (7), if $n$ is taken large enough.

But since, by \S.\,118, III, the numbers $S', R'$ are completely determined by the conditions (9), (10), and $S_{n-1}, R_{n-1}$ satisfy the same conditions by (6), (7), it follows that
\[
S' = S_{n-1}, \qquad R' = R_{n-1}.
\]

From this it now results further that the continued fraction with the closing number $x_n$
\weq{11}{(b_0, b_1, b_2 \ldots b_{m-1}, x_n)}
has the value
\[
\dfrac{R_n x_n + R_{n-1}}{S_n x_n + S_{n-1}},
\]
and therefore agrees with $y$.

If one now develops $x_n$ further into a continued fraction
\[
x_n = (a_n, a_{n+1}, a_{n+2} \ldots),
\]
one obtains from (2) and (11):
\[
x = (a_0, a_1, a_2 \ldots a_{n-1}, a_n, a_{n+1}, a_{n+2} \ldots)
\]
\[
y = (b_0, b_1, b_2 \ldots b_{m-1}, a_n, a_{n+1}, a_{n+2} \ldots),
\]
or, expressed in words, the theorem:

\gesperrt{VI. The continued fraction developments of two equivalent numbers agree with one another from a certain partial denominator on.}

We may add that, if the equivalence is a proper one ($\varepsilon = +1$), the number of partial denominators preceding the agreeing ones is even in both, or odd in both; and that, if the equivalence is improper, this number is even in the one and odd in the other.

That the theorem also holds conversely is easy to see. For if two continued fractions agree from a certain partial denominator on, they can be written in such a way that they have the same closing number. But the value of a continued fraction is always equivalent to each of its closing numbers [\S.\,117, (2)]; \gesperrt{hence two continued fractions with the same closing number are also equivalent to one another.}
